\documentclass[11pt, a4paper]{article}
\usepackage{geometry}
\usepackage[parfill]{parskip}
\usepackage{graphicx}
\usepackage{amsmath}
\usepackage{enumerate}

\title{Assignment 2 \\ IB3702 Mathematics for Machine learning}
\author{Harman Singh} % please overwrite this line with your name
\date{November 7, 2025} % please overwrite this line with the date of submission

\begin{document}
\maketitle

\pagebreak
\pagebreak
\section*{Problem 1}

\subsection*{exercise a}

$Av = \Lambda v$

First I am going to multiply $A$ to the first column of $Q$

So $v_1 = {[-a, 0, b]}^T$ with eigenvalue 25.

This gets me the equation:

\[
A \begin{pmatrix}
	-a \\
	0 \\
	b
\end{pmatrix} = 25 \begin{pmatrix}
	-a \\
	0 \\
	b
\end{pmatrix}
\]

Now I calculate $Av_1$ = 

\[
Av_1 = 
\begin{pmatrix}
 -2 & 0 & -36 \\
 0 & -3 & 0 \\
 -36 & 0 & -23	
\end{pmatrix} \begin{pmatrix}
	-a \\
	0 \\
	b
\end{pmatrix}
\]

Now we multiply horizontally on the left, with vertically on the right. So for example the element on the top left will be $(-2)(-a)$ and then $0(0)$ and then $-36(b)$ and so on.

\[
\begin{pmatrix}
	(-2)(-a) & + 0(0) & + (-36)(b) \\
	0(-a) & + (-3)(0) & + 0(b) \\
	(-36)(-a) & + 0(0) & + (-23)(b)
\end{pmatrix} = 
\begin{pmatrix}
	2a - 36b \\
	0 \\
	36a - 23b
\end{pmatrix}
\]

Now I equate this to $25v_1$:

\[
\begin{pmatrix}
	2a - 36b \\
	0 \\
	36a - 23b
\end{pmatrix} = 25 \begin{pmatrix}
	-a \\
	0 \\
	b
\end{pmatrix} = \begin{pmatrix}
	25a \\
	0 \\
	25b
\end{pmatrix}
\]

And now I can just solve the equations for $a$ and $b$:

\[
-2a - 36b = -25a \implies 27a - 36b = 0 \implies 27a = 36b \implies b = \frac{27}{36}a = \frac{3}{4}a
\]

\[
36a - 23b = 25b \implies 36a - 48b = 0 \implies 36a = 48b \implies b = \frac{36}{48}a = \frac{3}{4}a
\]

Since $Q$ orthogonally diagonalizes $A$, the following must be true
\[
	{(-a)}^2 + {(0)}^2 + {(b)}^2 = 1 \implies a^2 + b^2 = 1
\]

I can enter the value of $b$ I found above into this equation so we get:

\[
a^2 + {\left(\frac{3}{4}a\right)}^2 = 1 \implies a^2 + \frac{9}{16}a^2 = 1 \implies \frac{25}{16}a^2 = 1 \implies a^2 = \frac{16}{25} \implies a = \frac{4}{5}
\]

Then enter this value of $a$ into the equation for $b$:

\[b = \frac{3}{4}a =  \frac{3}{4} \cdot \frac{4}{5} = \frac{12}{20} = \frac{3}{5}\]


\[
\boxed{a = \frac{4}{5}, b = \frac{3}{5}}
\]



\subsection*{exercise b}
\ldots

\subsection*{exercise c}
\ldots

\subsection*{exercise d}
\ldots

\subsection*{exercise e}
\ldots




\pagebreak
\section*{Problem 2}

\subsection*{exercise a}
\ldots

\subsection*{exercise b}
\ldots

\subsection*{exercise c}
\ldots

\subsection*{exercise d}
\ldots

\subsection*{exercise e}
\ldots




\pagebreak
\section*{Problem 3}

% \begin{enumerate}[(a)]
% 	\item answer/solution to this problem
% 	\item answer/solution to this problem
% \end{enumerate}


\subsection*{exercise a}
\ldots

\subsection*{exercise b}
\ldots




\pagebreak
\section*{Problem 4}
\subsection*{exercise a}
\ldots

\subsection*{exercise b}
\ldots


\pagebreak
\section*{Problem 5}
\ldots




\ldots
\end{document}  